\documentclass[11pt]{article}
\usepackage[a4paper,margin=25mm]{geometry}
\usepackage{amsmath,amssymb,amsthm}
\newtheorem{lemma}{Lemma}\newtheorem{theorem}{Theorem}
\title{The Baer-subplane construction of Geometry 3.3}
\author{Private L100 component}\date{}
\begin{document}\maketitle
This working component concerns Geometry 3.3 on physical page 19 of the
retained 22-page Halle preprint of Gottschalk and Leemans. The journal
edition is not collated. The coordinate plane and its identification with
the retained 21-point permutation action are acquired from the accepted
projective-construction component. Its lines are literal five-point sets.

\begin{lemma}
A seven-point subset $B$ of $\operatorname{PG}(2,4)$ is a subplane of order
two if and only if every ambient line meets $B$ in one or three points.
\end{lemma}
\begin{proof}
Under the intersection condition, each pair of points of $B$ lies in a
unique three-point intersection. There are $\binom72/\binom32=7$ such
lines, with three through each point. Each meets the six others in its
three points, twice at each point: no two lines can share two points.
These seven points and seven lines therefore form a projective plane of
order two, embedded with its actual ambient joins.

Conversely, a subplane of order two has seven ambient lines meeting its
point set in three points; any other ambient line meets it in at most
one point, since two points determine a subplane line. Such another line
cannot avoid all seven subplane points. Otherwise its intersections with
the seven subplane lines would be distinct (two subplane lines meet in
the subplane), whereas an ambient line has only five points. Thus every
remaining line meets the subplane in one point.
\end{proof}

Enumerating all seven-subsets with this criterion is a complete finite
input construction, independent of the printed number of subplane orbits.
Orbits are obtained by closing each remaining subset under the exact
acquired point permutations. Every line pair is stored as two literal
five-point sets. Its image under a point permutation is the pair of the
images of those sets; no point/line label identification is assumed.

For an orbit $\mathcal B$, take the three types to be $\mathcal B$,
the unordered pairs of distinct lines, and the unordered pairs of distinct
points. Write $p_1p_2$ for the unique ambient line through $p_1,p_2$.
The three source incidences are exactly
\begin{align*}
B*\{\ell_1,\ell_2\}&\iff B\cap(\ell_1\cup\ell_2)=\ell_1\cap\ell_2,\\
B*\{p_1,p_2\}&\iff B\cap\{p_1,p_2\}=\varnothing\quad\text{and}\quad|B\cap p_1p_2|=3,\\
\{\ell_1,\ell_2\}*\{p_1,p_2\}&\iff
(\ell_1\cup\ell_2)\cap\{p_1,p_2\}=\varnothing\quad\text{and}\quad
\ell_1\cap\ell_2\cap p_1p_2\ne\varnothing.
\end{align*}
All are preserved by the acquired projective group.

\paragraph{Certificate contract (not yet executed).}
For each resulting subplane orbit, the finite certificate is to find a
type permutation and base objects fixed by the actual maximal parabolics
of acquired catalogue representative 10. An optional twist is literal
conjugation by a checked point normalizer. For that automorphism $\alpha$,
the candidate maps are $H_i g\mapsto x_i^{\alpha(g)}$.
The generator fixed-point tests prove well-definedness; the certificate
also checks every member of every right coset, object bijectivity, and
all incidences and nonincidences. A successful map transports the complete
accepted geometry, action and residue properties, including the actual
local-2T flags. Neither subgroup orders nor the printed diagram alone
can certify this join.

No execution or completion assertion is made by this working manuscript.
Other source constructions, butterflies, abstract group names and the
whole-source disposition remain separate.
\end{document}
